\documentclass[12pt,a4paper]{ctexart}

% ========== 页面与字体 ==========
\usepackage{geometry}
\geometry{left=2.6cm,right=2.6cm,top=2.5cm,bottom=2.5cm}

\usepackage{fontspec}
\setmainfont{Times New Roman}
\IfFontExistsTF{Consolas}
  {\setmonofont{Consolas}}
  {\setmonofont{Courier New}}

\usepackage{setspace}
\onehalfspacing

% ========== 数学与图表 ==========
\usepackage{amsmath, amsthm}
\usepackage{unicode-math}
\IfFontExistsTF{Cambria Math}{\setmathfont{Cambria Math}}{}
\usepackage{graphicx}
\usepackage{booktabs}
\usepackage{array}
\usepackage{caption}
\usepackage{subcaption}

% ========== 颜色与链接 ==========
\usepackage{xcolor}
\definecolor{mainblue}{HTML}{2F5597}
\definecolor{lightblue}{HTML}{EAF2FF}
\definecolor{darkgray}{HTML}{333333}

\usepackage[
    colorlinks=true,
    linkcolor=mainblue,
    citecolor=mainblue,
    urlcolor=mainblue
]{hyperref}

% ========== 标题样式 ==========
\usepackage{titlesec}

\titleformat{\section}
  {\Large\bfseries\color{mainblue}}
  {\thesection}{0.8em}{}

\titleformat{\subsection}
  {\large\bfseries\color{darkgray}}
  {\thesubsection}{0.8em}{}

\titleformat{\subsubsection}
  {\normalsize\bfseries}
  {\thesubsubsection}{0.8em}{}

% ========== 页眉页脚 ==========
\usepackage{fancyhdr}
\pagestyle{fancy}
\setlength{\headheight}{14.5pt}
\fancyhf{}
\fancyhead[L]{\color{mainblue}课程报告}
\fancyhead[R]{\color{mainblue}\leftmark}
\fancyfoot[C]{\thepage}

\renewcommand{\headrulewidth}{0.4pt}
\renewcommand{\footrulewidth}{0pt}

% ========== 代码样式 ==========
\usepackage{listings}

\lstset{
    basicstyle=\ttfamily\small,
    keywordstyle=\color{mainblue}\bfseries,
    commentstyle=\color{gray},
    stringstyle=\color{teal},
    numbers=left,
    numberstyle=\tiny\color{gray},
    frame=single,
    breaklines=true,
    backgroundcolor=\color{lightblue},
    rulecolor=\color{mainblue},
    tabsize=4
}

% ========== 漂亮的提示框 ==========
\usepackage[most]{tcolorbox}

\newtcolorbox{notebox}{
    colback=lightblue,
    colframe=mainblue,
    arc=2mm,
    boxrule=0.8pt,
    left=8pt,
    right=8pt,
    top=6pt,
    bottom=6pt,
    title=\textbf{提示}
}

% ========== 定理环境 ==========
\newtheorem{theorem}{定理}
\newtheorem{definition}{定义}
\newtheorem{example}{例}

% ========== 文档信息 ==========
\title{
    \vspace{2cm}
    \Huge \textbf{\color{mainblue}一份漂亮的 \LaTeX{} 中文模板} \\[0.5em]
    \Large Elegant Chinese LaTeX Template
}
\author{
    作者姓名 \\[0.3em]
    \small 学校 / 单位
}
\date{\today}

% ========== 正文开始 ==========
\begin{document}

\maketitle
\thispagestyle{empty}

\vspace{2cm}

\begin{abstract}
本文档提供了一个简洁、美观、适合中文写作的 \LaTeX{} 模板。它包含标题页、目录、页眉页脚、数学公式、表格、代码块、提示框和定理环境等常用结构。
\end{abstract}

\newpage
\tableofcontents
\newpage

\section{引言}

这是一个适合课程论文、实验报告、读书笔记和技术文档的中文 \LaTeX{} 模板。

\begin{notebox}
这个模板推荐使用 \textbf{XeLaTeX} 编译，能够较好地支持中文字体。
\end{notebox}

\section{数学公式示例}

下面是一个行内公式：$E = mc^2$。

下面是一个独立公式：

\[
    \int_{-\infty}^{+\infty} e^{-x^2} \, dx = \sqrt{\pi}
\]

也可以写带编号的公式：

\begin{equation}
    a^2 + b^2 = c^2
\end{equation}

\section{表格示例}

\begin{table}[htbp]
\centering
\caption{实验结果示例}
\begin{tabular}{cccc}
\toprule
编号 & 方法 & 准确率 & 时间/s \\
\midrule
1 & 方法 A & 92.5\% & 1.23 \\
2 & 方法 B & 94.1\% & 1.56 \\
3 & 方法 C & 95.3\% & 1.48 \\
\bottomrule
\end{tabular}
\end{table}

\section{图片示例}

\begin{figure}[htbp]
\centering
% \includegraphics[width=0.65\textwidth]{example-image}
\caption{这里放置图片标题}
\end{figure}

\section{代码示例}

\begin{lstlisting}[language=Python, caption={Python 示例代码}]
def hello(name):
    print(f"Hello, {name}!")

hello("LaTeX")
\end{lstlisting}

\section{定理环境示例}

\begin{definition}
若函数 $f(x)$ 在点 $x_0$ 的某个邻域内有定义，并且极限
\[
    \lim_{x \to x_0} f(x)
\]
存在，则称该极限为函数在 $x_0$ 处的极限。
\end{definition}

\begin{theorem}
若函数 $f(x)$ 在闭区间 $[a,b]$ 上连续，则 $f(x)$ 在该区间上有界。
\end{theorem}

\section{结论}

本文档展示了一个基础但美观的中文 \LaTeX{} 模板。你可以在此基础上继续添加参考文献、算法环境、封面图片或学校格式要求。

\end{document}
